\documentclass[11pt]{article}
\usepackage[utf8]{inputenc}
\usepackage{amsmath, amssymb, amsthm}
\usepackage[margin=1in]{geometry}
\usepackage{hyperref}
\usepackage{listings}
\usepackage{color}

\definecolor{codegray}{rgb}{0.5,0.5,0.5}
\definecolor{codepurple}{rgb}{0.58,0,0.82}
\definecolor{backcolour}{rgb}{0.98,0.98,0.98}

\lstset{ 
    backgroundcolor=\color{backcolour},   
    basicstyle=\ttfamily\footnotesize,
    breakatwhitespace=false,         
    breaklines=true,                 
    captionpos=b,                    
    keepspaces=true,                 
    numbers=none,                    
    showspaces=false,                
    showstringspaces=false,
    showtabs=false,                  
    tabsize=2,
    morekeywords={def, noncomputable, Prop, let, in, fun, if, then, else, match, with, where},
    mathescape=true
}

\title{Definitions in \texttt{Sarason/} and \texttt{R2.lean}}
\author{Lean Formalization Project}
\date{\today}

\lstset{
  literate={×}{{$\times$}}1 {γ}{{$\gamma$}}1 {δ}{{$\delta$}}1 {ε}{{$\varepsilon$}}1 {‖}{{$\|$}}1 {₀}{{$_0$}}1 {ℂ}{{$\mathbb{C}$}}1 {ℕ}{{$\mathbb{N}$}}1 {ℝ}{{$\mathbb{R}$}}1 {→}{{$\to$}}1 {∀}{{$\forall$}}1 {∃}{{$\exists$}}1 {∧}{{$\wedge$}}1 {≤}{{$\le$}}1 {≥}{{$\ge$}}1 {▸}{{$\triangleright$}}1 {⟨}{{$\langle$}}1 {⟩}{{$\rangle$}}1 {≠}{{$\neq$}}1
}
\begin{document}
\maketitle

This document provides a comprehensive list of all definitions manually formalized in the \texttt{ComplexAnalysis} project, specifically within \texttt{R2.lean} and the \texttt{Sarason/} directory. For each definition, the raw Lean code is presented first, followed by its mathematical interpretation.

\section{\texttt{R2.lean}}

\begin{itemize}
    \item \textbf{\texttt{sqNorm}} ($x \in \mathbb{R}^2 \to \mathbb{R}$): The squared Euclidean norm.
\begin{lstlisting}[language=Caml, mathescape=true]
noncomputable def sqNorm (x : $\mathbb{R}$ $\times$ $\mathbb{R}$) : $\mathbb{R}$ := x.1^2 + x.2^2
\end{lstlisting}
    \[ \text{sqNorm}(x) = x_1^2 + x_2^2 \]

    \item \textbf{\texttt{sqDist}} ($x, y \in \mathbb{R}^2 \to \mathbb{R}$): The squared Euclidean distance.
\begin{lstlisting}[language=Caml, mathescape=true]
noncomputable def sqDist (x y : $\mathbb{R}$ $\times$ $\mathbb{R}$) : $\mathbb{R}$ :=
  (x.1 - y.1)^2 + (x.2 - y.2)^2
\end{lstlisting}
    \[ \text{sqDist}(x, y) = \text{sqNorm}(x - y) \]

    \item \textbf{\texttt{euclideanNorm}} ($x \in \mathbb{R}^2 \to \mathbb{R}$): The standard Euclidean norm.
\begin{lstlisting}[language=Caml, mathescape=true]
noncomputable def euclideanNorm (x : $\mathbb{R}$ $\times$ $\mathbb{R}$) : $\mathbb{R}$ :=
  Real.sqrt (sqNorm x)
\end{lstlisting}
    \[ \text{euclideanNorm}(x) = \sqrt{x_1^2 + x_2^2} \]

    \item \textbf{\texttt{euclideanDist}} ($x, y \in \mathbb{R}^2 \to \mathbb{R}$): The Euclidean distance.
\begin{lstlisting}[language=Caml, mathescape=true]
noncomputable def euclideanDist (x y : $\mathbb{R}$ $\times$ $\mathbb{R}$) : $\mathbb{R}$ :=
  Real.sqrt (sqDist x y)
\end{lstlisting}
    \[ \text{euclideanDist}(x, y) = \sqrt{\text{sqDist}(x, y)} \]

    \item \textbf{\texttt{LimitRtoR2}} ($f : \mathbb{R} \to \mathbb{R}^2$, $a \in \mathbb{R}$, $L \in \mathbb{R}^2$):
\begin{lstlisting}[language=Caml, mathescape=true]
def LimitRtoR2 (f : $\mathbb{R}$ $\to$ $\mathbb{R}$ $\times$ $\mathbb{R}$) (a : $\mathbb{R}$) (L : $\mathbb{R}$ $\times$ $\mathbb{R}$) : Prop :=
  $\forall$ $\varepsilon$ > 0, $\exists$ $\delta$ > 0, $\forall$ x : $\mathbb{R}$, 0 < |x - a| $\wedge$ |x - a| < $\delta$ $\to$ euclideanDist (f x) L < $\varepsilon$
\end{lstlisting}
    \[ \lim_{x \to a} f(x) = L \iff \forall \epsilon > 0, \exists \delta > 0, \forall x, 0 < |x - a| < \delta \implies \text{euclideanDist}(f(x), L) < \epsilon \]

    \item \textbf{\texttt{LimitR2toR}} ($f : \mathbb{R}^2 \to \mathbb{R}$, $a \in \mathbb{R}^2$, $L \in \mathbb{R}$):
\begin{lstlisting}[language=Caml, mathescape=true]
def LimitR2toR (f : $\mathbb{R}$ $\times$ $\mathbb{R}$ $\to$ $\mathbb{R}$) (a : $\mathbb{R}$ $\times$ $\mathbb{R}$) (L : $\mathbb{R}$) : Prop :=
  $\forall$ $\varepsilon$ > 0, $\exists$ $\delta$ > 0, $\forall$ x : $\mathbb{R}$ $\times$ $\mathbb{R}$, 0 < euclideanDist x a $\wedge$ euclideanDist x a < $\delta$ $\to$ |f x - L| < $\varepsilon$
\end{lstlisting}
    \[ \lim_{x \to a} f(x) = L \iff \forall \epsilon > 0, \exists \delta > 0, \forall x \in \mathbb{R}^2, 0 < \text{euclideanDist}(x, a) < \delta \implies |f(x) - L| < \epsilon \]

    \item \textbf{\texttt{LimitR2toR2}} ($f : \mathbb{R}^2 \to \mathbb{R}^2$, $a \in \mathbb{R}^2$, $L \in \mathbb{R}^2$):
\begin{lstlisting}[language=Caml, mathescape=true]
def LimitR2toR2 (f : $\mathbb{R}$ $\times$ $\mathbb{R}$ $\to$ $\mathbb{R}$ $\times$ $\mathbb{R}$) (a : $\mathbb{R}$ $\times$ $\mathbb{R}$) (L : $\mathbb{R}$ $\times$ $\mathbb{R}$) : Prop :=
  $\forall$ $\varepsilon$ > 0, $\exists$ $\delta$ > 0, $\forall$ x : $\mathbb{R}$ $\times$ $\mathbb{R}$, 0 < euclideanDist x a $\wedge$ euclideanDist x a < $\delta$ $\to$ euclideanDist (f x) L < $\varepsilon$
\end{lstlisting}
    \[ \lim_{x \to a} f(x) = L \iff \forall \epsilon > 0, \exists \delta > 0, \forall x \in \mathbb{R}^2, 0 < \text{euclideanDist}(x, a) < \delta \implies \text{euclideanDist}(f(x), L) < \epsilon \]

    \item \textbf{\texttt{ConvergesR2}} ($\text{seq} : \mathbb{N} \to \mathbb{R}^2$, $L \in \mathbb{R}^2$):
\begin{lstlisting}[language=Caml, mathescape=true]
def ConvergesR2 (seq : $\mathbb{N}$ $\to$ $\mathbb{R}$ $\times$ $\mathbb{R}$) (L : $\mathbb{R}$ $\times$ $\mathbb{R}$): Prop :=
  $\forall$ $\varepsilon$ > 0, $\exists$ N : $\mathbb{N}$, $\forall$ n $\ge$ N, euclideanDist (seq n) L < $\varepsilon$
\end{lstlisting}
    \[ \lim_{n \to \infty} \text{seq}(n) = L \iff \forall \epsilon > 0, \exists N \in \mathbb{N}, \forall n \ge N, \text{euclideanDist}(\text{seq}(n), L) < \epsilon \]

    \item \textbf{\texttt{HasFDerivAt\_R2\_eps}} ($u, v : \mathbb{R}^2 \to \mathbb{R}$, $u_{x0}, u_{y0}, v_{x0}, v_{y0} \in \mathbb{R}$, $a \in \mathbb{R}^2$): Defines the Fréchet derivative of $F(x,y) = (u(x,y), v(x,y))$ at $a$. Let $L(h) = \begin{pmatrix} u_{x0} & u_{y0} \\ v_{x0} & v_{y0} \end{pmatrix} h$.
\begin{lstlisting}[language=Caml, mathescape=true]
def HasFDerivAt_R2_eps (u v : $\mathbb{R}$ $\times$ $\mathbb{R}$ $\to$ $\mathbb{R}$) (ux$_0$ uy$_0$ vx$_0$ vy$_0$ : $\mathbb{R}$) (a : $\mathbb{R}$ $\times$ $\mathbb{R}$) : Prop :=
  LimitR2toR (fun h => 
    euclideanNorm (
      (u (a.1 + h.1, a.2 + h.2) - u a) - (ux$_0$ * h.1 + uy$_0$ * h.2), 
      (v (a.1 + h.1, a.2 + h.2) - v a) - (vx$_0$ * h.1 + vy$_0$ * h.2)
    ) / euclideanNorm h
  ) (0, 0) 0
\end{lstlisting}
    \[ \lim_{h \to 0} \frac{\left\| F(a+h) - F(a) - \begin{pmatrix} u_{x0} & u_{y0} \\ v_{x0} & v_{y0} \end{pmatrix} h \right\|}{\|h\|} = 0 \]
    \[ \lim_{(h_1, h_2) \to (0,0)} \frac{\left\| \begin{pmatrix} u(a_1+h_1, a_2+h_2) - u(a_1, a_2) \\ v(a_1+h_1, a_2+h_2) - v(a_1, a_2) \end{pmatrix} - \begin{pmatrix} u_{x0} & u_{y0} \\ v_{x0} & v_{y0} \end{pmatrix} \begin{pmatrix} h_1 \\ h_2 \end{pmatrix} \right\|}{\| (h_1, h_2) \|} = 0 \]


    \item \textbf{\texttt{DifferentiableAt\_R2\_eps}} ($u, v : \mathbb{R}^2 \to \mathbb{R}$, $a \in \mathbb{R}^2$):
\begin{lstlisting}[language=Caml, mathescape=true]
def DifferentiableAt_R2_eps (u v : $\mathbb{R}$ $\times$ $\mathbb{R}$ $\to$ $\mathbb{R}$) (a : $\mathbb{R}$ $\times$ $\mathbb{R}$) : Prop :=
  $\exists$ (ux$_0$ uy$_0$ vx$_0$ vy$_0$ : $\mathbb{R}$), HasFDerivAt_R2_eps u v ux$_0$ uy$_0$ vx$_0$ vy$_0$ a
\end{lstlisting}
    \[ \exists u_{x0}, u_{y0}, v_{x0}, v_{y0} \in \mathbb{R}, \text{HasFDerivAt\_R2\_eps}(u, v, u_{x0}, u_{y0}, v_{x0}, v_{y0}, a) \]

    \item \textbf{\texttt{HasDerivAt\_RtoR2\_eps}} ($f : \mathbb{R} \to \mathbb{R}^2$, $f'_0 \in \mathbb{R}^2$, $t_0 \in \mathbb{R}$): Defines the derivative of a parametric curve in $\mathbb{R}^2$.
\begin{lstlisting}[language=Caml, mathescape=true]
def HasDerivAt_RtoR2_eps (f : $\mathbb{R}$ $\to$ $\mathbb{R}$ $\times$ $\mathbb{R}$) (f'$_0$ : $\mathbb{R}$ $\times$ $\mathbb{R}$) (t$_0$ : $\mathbb{R}$) : Prop :=
  $\forall$ $\varepsilon$ > 0, $\exists$ $\delta$ > 0, $\forall$ t : $\mathbb{R}$, 0 < |t - t$_0$| $\wedge$ |t - t$_0$| < $\delta$ $\to$
    euclideanDist (f t) (f t$_0$ + (f'$_0$.1 * (t - t$_0$), f'$_0$.2 * (t - t$_0$))) < $\varepsilon$ * |t - t$_0$|
\end{lstlisting}
    \[ \lim_{t \to t_0} (f(t) - f(t_0) - (t - t_0)f'_0) = 0 \]

    \item \textbf{\texttt{HasPartialDerivX\_R2\_eps}} ($u : \mathbb{R}^2 \to \mathbb{R}$, $ux_0 \in \mathbb{R}$, $p_0 \in \mathbb{R}^2$):
\begin{lstlisting}[language=Caml, mathescape=true]
def HasPartialDerivX_R2_eps (u : $\mathbb{R}$ $\times$ $\mathbb{R}$ $\to$ $\mathbb{R}$) (ux$_0$ : $\mathbb{R}$) (p$_0$ : $\mathbb{R}$ $\times$ $\mathbb{R}$) : Prop :=
  $\forall$ $\varepsilon$ > 0, $\exists$ $\delta$ > 0, $\forall$ x : $\mathbb{R}$, 0 < |x - p$_0$.1| $\wedge$ |x - p$_0$.1| < $\delta$ $\to$
    |(u (x, p$_0$.2) - u p$_0$) / (x - p$_0$.1) - ux$_0$| < $\varepsilon$
\end{lstlisting}
    \[ \lim_{x \to p_{0,x}} \frac{u(x, p_{0,y}) - u(p_0)}{x - p_{0,x}} = ux_0 \]

    \item \textbf{\texttt{HasPartialDerivY\_R2\_eps}} ($u : \mathbb{R}^2 \to \mathbb{R}$, $uy_0 \in \mathbb{R}$, $p_0 \in \mathbb{R}^2$):
\begin{lstlisting}[language=Caml, mathescape=true]
def HasPartialDerivY_R2_eps (u : $\mathbb{R}$ $\times$ $\mathbb{R}$ $\to$ $\mathbb{R}$) (uy$_0$ : $\mathbb{R}$) (p$_0$ : $\mathbb{R}$ $\times$ $\mathbb{R}$) : Prop :=
  $\forall$ $\varepsilon$ > 0, $\exists$ $\delta$ > 0, $\forall$ y : $\mathbb{R}$, 0 < |y - p$_0$.2| $\wedge$ |y - p$_0$.2| < $\delta$ $\to$
    |(u (p$_0$.1, y) - u p$_0$) / (y - p$_0$.2) - uy$_0$| < $\varepsilon$

/-- A real-valued function of two real variables is harmonic on an open set G
if its first and second partial derivatives are continuous and satisfy Laplace's equation. -/
\end{lstlisting}
    \[ \lim_{y \to p_{0,y}} \frac{u(p_{0,x}, y) - u(p_0)}{y - p_{0,y}} = uy_0 \]

    \item \textbf{\texttt{TwiceContinuouslyDifferentiable\_R2\_eps}} ($u : \mathbb{R}^2 \to \mathbb{R}$, $G \subset \mathbb{R}^2$): Defines twice continuous differentiability ($C^2$) for a real-valued function on an open subset $G \subset \mathbb{R}^2$.
\begin{lstlisting}[language=Caml, mathescape=true]
def TwiceContinuouslyDifferentiable_R2_eps (u : $\mathbb{R}$ $\times$ $\mathbb{R}$ $\to$ $\mathbb{R}$) (G : Set ($\mathbb{R}$ $\times$ $\mathbb{R}$)) : Prop :=
  IsOpen G $\wedge$
  $\exists$ (ux uy uxx uxy uyx uyy : $\mathbb{R}$ $\times$ $\mathbb{R}$ $\to$ $\mathbb{R}$),
    ($\forall$ p $\in$ G, HasPartialDerivX_R2_eps u (ux p) p) $\wedge$
    ($\forall$ p $\in$ G, HasPartialDerivY_R2_eps u (uy p) p) $\wedge$
    (ContinuousOn ux G) $\wedge$
    (ContinuousOn uy G) $\wedge$
    ($\forall$ p $\in$ G, HasPartialDerivX_R2_eps ux (uxx p) p) $\wedge$
    ($\forall$ p $\in$ G, HasPartialDerivY_R2_eps ux (uxy p) p) $\wedge$
    ($\forall$ p $\in$ G, HasPartialDerivX_R2_eps uy (uyx p) p) $\wedge$
    ($\forall$ p $\in$ G, HasPartialDerivY_R2_eps uy (uyy p) p) $\wedge$
    (ContinuousOn uxx G) $\wedge$
    (ContinuousOn uxy G) $\wedge$
    (ContinuousOn uyx G) $\wedge$
    (ContinuousOn uyy G)
\end{lstlisting}
    \[ u \in C^2(G, \mathbb{R}) \iff u_x, u_y, u_{xx}, u_{xy}, u_{yx}, u_{yy} \text{ exist and are continuous on the open set } G \]

    \item \textbf{\texttt{Harmonic\_R2\_eps}} ($u : \mathbb{R}^2 \to \mathbb{R}$, $G \subset \mathbb{R}^2$):
\begin{lstlisting}[language=Caml, mathescape=true]
def Harmonic_R2_eps (u : $\mathbb{R}$ $\times$ $\mathbb{R}$ $\to$ $\mathbb{R}$) (G : Set ($\mathbb{R}$ $\times$ $\mathbb{R}$)) : Prop :=
  IsOpen G $\wedge$
  $\exists$ (ux uy uxx uxy uyx uyy : $\mathbb{R}$ $\times$ $\mathbb{R}$ $\to$ $\mathbb{R}$),
    ($\forall$ p $\in$ G, HasPartialDerivX_R2_eps u (ux p) p) $\wedge$
    ($\forall$ p $\in$ G, HasPartialDerivY_R2_eps u (uy p) p) $\wedge$
    (ContinuousOn ux G) $\wedge$
    (ContinuousOn uy G) $\wedge$
    ($\forall$ p $\in$ G, HasPartialDerivX_R2_eps ux (uxx p) p) $\wedge$
    ($\forall$ p $\in$ G, HasPartialDerivY_R2_eps ux (uxy p) p) $\wedge$
    ($\forall$ p $\in$ G, HasPartialDerivX_R2_eps uy (uyx p) p) $\wedge$
    ($\forall$ p $\in$ G, HasPartialDerivY_R2_eps uy (uyy p) p) $\wedge$
    (ContinuousOn uxx G) $\wedge$
    (ContinuousOn uxy G) $\wedge$
    (ContinuousOn uyx G) $\wedge$
    (ContinuousOn uyy G) $\wedge$
    ($\forall$ p $\in$ G, uxx p + uyy p = 0)
\end{lstlisting}
    \[ \text{First and second partials continuous on } G, \quad \Delta u = u_{xx} + u_{yy} = 0 \]
\end{itemize}


\section{\texttt{Sarason/Definitions.lean}}

\begin{itemize}
    \item \textbf{\texttt{HasLimitAt\_eps}} ($f : \mathbb{C} \to \mathbb{C}$, $L \in \mathbb{C}$, $z_0 \in \mathbb{C}$):
\begin{lstlisting}[language=Caml, mathescape=true]
def HasLimitAt_eps (f : $\mathbb{C}$ $\to$ $\mathbb{C}$) (L : $\mathbb{C}$) (z$_0$ : $\mathbb{C}$) : Prop :=
  $\forall$ $\varepsilon$ > 0, $\exists$ $\delta$ > 0, $\forall$ z, 0 < $\|$z - z$_0$$\|$ $\wedge$ $\|$z - z$_0$$\|$ < $\delta$ $\to$ $\|$f z - L$\|$ < $\varepsilon$

--  Weierstrass definition of a limit at infinity.
\end{lstlisting}
    \[ \lim_{z \to z_0} f(z) = L \]

    \item \textbf{\texttt{HasLimitAtInf\_eps}} ($f : \mathbb{C} \to \mathbb{R}$, $L \in \mathbb{R}$):
\begin{lstlisting}[language=Caml, mathescape=true]
def HasLimitAtInf_eps (f : $\mathbb{C}$ $\to$ $\mathbb{R}$) (L : $\mathbb{R}$) : Prop :=
  $\forall$ $\varepsilon$ > 0, $\exists$ M > 0, $\forall$ z, $\|$z$\|$ > M $\to$ |f z - L| < $\varepsilon$

--  Weierstrass definition of having a complex derivative at z$_0$.
\end{lstlisting}
    \[ \lim_{|z| \to \infty} f(z) = L \]

    \item \textbf{\texttt{HasDerivAt\_eps}} ($f : \mathbb{C} \to \mathbb{C}$, $f'_0 \in \mathbb{C}$, $z_0 \in \mathbb{C}$): Complex differentiability at a point.
\begin{lstlisting}[language=Caml, mathescape=true]
def HasDerivAt_eps (f : $\mathbb{C}$ $\to$ $\mathbb{C}$) (f'$_0$ : $\mathbb{C}$) (z$_0$ : $\mathbb{C}$) : Prop :=
  HasLimitAt_eps (fun z => (f z - f z$_0$) / (z - z$_0$)) f'$_0$ z$_0$

--  Weierstrass definition of complex differentiability at z$_0$.
\end{lstlisting}
    \[ \lim_{z \to z_0} \frac{f(z) - f(z_0)}{z - z_0} = f'_0 \]

    \item \textbf{\texttt{DifferentiableAt\_eps}} ($f : \mathbb{C} \to \mathbb{C}$, $z_0 \in \mathbb{C}$):
\begin{lstlisting}[language=Caml, mathescape=true]
def DifferentiableAt_eps (f : $\mathbb{C}$ $\to$ $\mathbb{C}$) (z$_0$ : $\mathbb{C}$) : Prop :=
  $\exists$ f'$_0$, HasDerivAt_eps f f'$_0$ z$_0$

-- Weierstrass definition of partial derivative with respect to x.
\end{lstlisting}
    \[ f \text{ is differentiable at } z_0 \]

    \item \textbf{\texttt{HasPartialDerivX\_C\_to\_R\_eps}} / \textbf{\texttt{HasPartialDerivY\_C\_to\_R\_eps}}: Partial derivatives of a real-valued function of a complex variable (identifying $\mathbb{C} \cong \mathbb{R}^2$).
\begin{lstlisting}[language=Caml, mathescape=true]
def HasPartialDerivX_C_to_R_eps (u : $\mathbb{C}$ $\to$ $\mathbb{R}$) (ux$_0$ : $\mathbb{R}$) (z$_0$ : $\mathbb{C}$) : Prop :=
  $\forall$ $\varepsilon$ > 0, $\exists$ $\delta$ > 0, $\forall$ x : $\mathbb{R}$, 0 < |x - z$_0$.re| $\wedge$ |x - z$_0$.re| < $\delta$ $\to$
    |(u ((x : $\mathbb{C}$) + (z$_0$.im : $\mathbb{C}$) * I) - u z$_0$) / (x - z$_0$.re) - ux$_0$| < $\varepsilon$

-- Weierstrass definition of partial derivative with respect to y.
\end{lstlisting}
    \[ \frac{\partial u}{\partial x}(z_0) = u_{x0} \]

    \item \textbf{\texttt{HasPartialDerivY\_C\_to\_R\_eps}}
\begin{lstlisting}[language=Caml, mathescape=true]
def HasPartialDerivY_C_to_R_eps (u : $\mathbb{C}$ $\to$ $\mathbb{R}$) (uy$_0$ : $\mathbb{R}$) (z$_0$ : $\mathbb{C}$) : Prop :=
  $\forall$ $\varepsilon$ > 0, $\exists$ $\delta$ > 0, $\forall$ y : $\mathbb{R}$, 0 < |y - z$_0$.im| $\wedge$ |y - z$_0$.im| < $\delta$ $\to$
    |(u ((z$_0$.re : $\mathbb{C}$) + (y : $\mathbb{C}$) * I) - u z$_0$) / (y - z$_0$.im) - uy$_0$| < $\varepsilon$

--  Equivalence theorem between Mathlib's limit at a point (using punctured neighborhoods) and Weierstrass epsilon-delta limit.
\end{lstlisting}
    \[ \frac{\partial u}{\partial y}(z_0) = u_{y0} \]

    \item \textbf{\texttt{deriv\_eps}} ($f : \mathbb{C} \to \mathbb{C}$, $z_0 \in \mathbb{C}$, $h : \texttt{DifferentiableAt\_eps}\ f\ z_0$): The classical derivative value $f'(z_0)$ extracted via choice.
\begin{lstlisting}[language=Caml, mathescape=true]
noncomputable def deriv_eps (f : $\mathbb{C}$ $\to$ $\mathbb{C}$) (z$_0$ : $\mathbb{C}$) (h : DifferentiableAt_eps f z$_0$) : $\mathbb{C}$ :=
  Classical.choose h
\end{lstlisting}
    \[ \text{deriv\_eps}(f, z_0) = f'(z_0) \]

    \item \textbf{\texttt{HasDerivAt\_R\_to\_C\_eps}} ($f : \mathbb{R} \to \mathbb{C}$, $f'_0 \in \mathbb{C}$, $t_0 \in \mathbb{R}$): Derivative of a complex-valued path parameterized by a real variable.
\begin{lstlisting}[language=Caml, mathescape=true]
def HasDerivAt_R_to_C_eps (f : $\mathbb{R}$ $\to$ $\mathbb{C}$) (f'$_0$ : $\mathbb{C}$) (t$_0$ : $\mathbb{R}$) : Prop :=
  $\forall$ $\varepsilon$ > 0, $\exists$ $\delta$ > 0, $\forall$ t : $\mathbb{R}$, 0 < |t - t$_0$| $\wedge$ |t - t$_0$| < $\delta$ $\to$ $\|$(f t - f t$_0$) / ((t - t$_0$ : $\mathbb{R}$) : $\mathbb{C}$) - f'$_0$$\|$ < $\varepsilon$
\end{lstlisting}
    \[ \lim_{t \to t_0} \frac{f(t) - f(t_0)}{t - t_0} = f'_0 \]

    \item \textbf{\texttt{DifferentiableAt\_R\_to\_C\_eps}} ($f : \mathbb{R} \to \mathbb{C}$, $t_0 \in \mathbb{R}$):
\begin{lstlisting}[language=Caml, mathescape=true]
def DifferentiableAt_R_to_C_eps (f : $\mathbb{R}$ $\to$ $\mathbb{C}$) (t$_0$ : $\mathbb{R}$) : Prop :=
  $\exists$ f'$_0$, HasDerivAt_R_to_C_eps f f'$_0$ t$_0$

/-- The classical derivative value extracted via choice, or 0 if not differentiable. -/
\end{lstlisting}
    \[ \exists f'_0 \in \mathbb{C}, \text{HasDerivAt\_R\_to\_C\_eps}(f, f'_0, t_0) \]

    \item \textbf{\texttt{deriv\_R\_to\_C\_eps}} ($f : \mathbb{R} \to \mathbb{C}$, $t_0 \in \mathbb{R}$):
\begin{lstlisting}[language=Caml, mathescape=true]
noncomputable def deriv_R_to_C_eps (f : $\mathbb{R}$ $\to$ $\mathbb{C}$) (t$_0$ : $\mathbb{R}$) : $\mathbb{C}$ :=
  if h : DifferentiableAt_R_to_C_eps f t$_0$ then Classical.choose h else 0
\end{lstlisting}
    \[ \text{deriv\_R\_to\_C\_eps}(f, t_0) = \begin{cases} f'(t_0) & \text{if differentiable} \\ 0 & \text{otherwise} \end{cases} \]

    \item \textbf{\texttt{HolomorphicAt\_eps}} ($f : \mathbb{C} \to \mathbb{C}$, $z_0 \in \mathbb{C}$): $f$ is complex-differentiable in some open neighborhood around $z_0$.
\begin{lstlisting}[language=Caml, mathescape=true]
def HolomorphicAt_eps (f : $\mathbb{C}$ $\to$ $\mathbb{C}$) (z$_0$ : $\mathbb{C}$) : Prop :=
  DifferentiableAt_eps f z$_0$

/-- A function is holomorphic on a set if it is complex differentiable at every point of the set. -/
\end{lstlisting}
    \[ f \text{ is holomorphic at } z_0 \iff f \text{ is complex-differentiable on some open ball around } z_0 \]

    \item \textbf{\texttt{HolomorphicOn\_eps}} ($f : \mathbb{C} \to \mathbb{C}$, $G \subset \mathbb{C}$): $f$ is holomorphic at every point in the open set $G$.
\begin{lstlisting}[language=Caml, mathescape=true]
def HolomorphicOn_eps (f : $\mathbb{C}$ $\to$ $\mathbb{C}$) (G : Set $\mathbb{C}$) : Prop :=
  $\forall$ z $\in$ G, DifferentiableAt_eps f z
\end{lstlisting}
    \[ f \text{ is holomorphic on } G \]

    \item \textbf{\texttt{polar\_partial\_r}} ($ux, uy : \mathbb{R}^2 \to \mathbb{R}$, $p \in \mathbb{R}^2$): Directional derivative in the radial direction for polar coordinates.
\begin{lstlisting}[language=Caml, mathescape=true]
noncomputable def polar_partial_r (ux uy : $\mathbb{R}$ $\times$ $\mathbb{R}$ $\to$ $\mathbb{R}$) (p : $\mathbb{R}$ $\times$ $\mathbb{R}$) : $\mathbb{R}$ :=
  ux (p.1 * Real.cos p.2, p.1 * Real.sin p.2) * Real.cos p.2 +
  uy (p.1 * Real.cos p.2, p.1 * Real.sin p.2) * Real.sin p.2
\end{lstlisting}
    \[ U_r(r, \theta) = u_x(r \cos \theta, r \sin \theta) \cos \theta + u_y(r \cos \theta, r \sin \theta) \sin \theta \]

    \item \textbf{\texttt{polar\_partial\_theta}} ($ux, uy : \mathbb{R}^2 \to \mathbb{R}$, $p \in \mathbb{R}^2$): Directional derivative in the angular direction for polar coordinates.
\begin{lstlisting}[language=Caml, mathescape=true]
noncomputable def polar_partial_theta (ux uy : $\mathbb{R}$ $\times$ $\mathbb{R}$ $\to$ $\mathbb{R}$) (p : $\mathbb{R}$ $\times$ $\mathbb{R}$) : $\mathbb{R}$ :=
  -ux (p.1 * Real.cos p.2, p.1 * Real.sin p.2) * p.1 * Real.sin p.2 +
  uy (p.1 * Real.cos p.2, p.1 * Real.sin p.2) * p.1 * Real.cos p.2
\end{lstlisting}
    \[ U_\theta(r, \theta) = -r \sin \theta \ u_x(r \cos \theta, r \sin \theta) + r \cos \theta \ u_y(r \cos \theta, r \sin \theta) \]

    \item \textbf{\texttt{polar\_partial\_rr}} ($uxx, uxy, uyy : \mathbb{R}^2 \to \mathbb{R}$, $p \in \mathbb{R}^2$): Second partial derivative with respect to $r$ in polar coordinates.
\begin{lstlisting}[language=Caml, mathescape=true]
noncomputable def polar_partial_rr (uxx uxy uyy : $\mathbb{R}$ $\times$ $\mathbb{R}$ $\to$ $\mathbb{R}$) (p : $\mathbb{R}$ $\times$ $\mathbb{R}$) : $\mathbb{R}$ :=
  Real.cos p.2 * polar_partial_r uxx uxy p + Real.sin p.2 * polar_partial_r uxy uyy p
\end{lstlisting}
    \[ U_{rr}(r, \theta) = \cos \theta \, (u_{xx})_r + \sin \theta \, (u_{xy})_r = \cos^2 \theta \, u_{xx} + 2 \sin \theta \cos \theta \, u_{xy} + \sin^2 \theta \, u_{yy} \]

    \item \textbf{\texttt{polar\_partial\_thetatheta}} ($ux, uy, uxx, uxy, uyy : \mathbb{R}^2 \to \mathbb{R}$, $p \in \mathbb{R}^2$): Second partial derivative with respect to $\theta$ in polar coordinates.
\begin{lstlisting}[language=Caml, mathescape=true]
noncomputable def polar_partial_thetatheta (ux uy uxx uxy uyy : $\mathbb{R}$ $\times$ $\mathbb{R}$ $\to$ $\mathbb{R}$) (p : $\mathbb{R}$ $\times$ $\mathbb{R}$) : $\mathbb{R}$ :=
  -p.1 * Real.cos p.2 * ux (p.1 * Real.cos p.2, p.1 * Real.sin p.2)
  - p.1 * Real.sin p.2 * polar_partial_theta uxx uxy p
  - p.1 * Real.sin p.2 * uy (p.1 * Real.cos p.2, p.1 * Real.sin p.2)
  + p.1 * Real.cos p.2 * polar_partial_theta uxy uyy p
\end{lstlisting}
    \[ U_{\theta\theta}(r, \theta) = -r(u_x \cos\theta + u_y \sin\theta) - r\sin\theta\,(u_{xx})_\theta + r\cos\theta\,(u_{xy})_\theta \]

    \item \textbf{\texttt{TwiceContinuouslyDifferentiable\_C\_eps}} ($f : \mathbb{C} \to \mathbb{C}$, $G \subset \mathbb{C}$): $f$ is twice continuously differentiable on the open set $G$ if its real and imaginary parts are.
\begin{lstlisting}[language=Caml, mathescape=true]
def TwiceContinuouslyDifferentiable_C_eps (f : $\mathbb{C}$ $\to$ $\mathbb{C}$) (G : Set $\mathbb{C}$) : Prop :=
  ComplexAnalysis.R2.TwiceContinuouslyDifferentiable_R2_eps (fun p => (f (p.1 + p.2 * I)).re) {p : $\mathbb{R}$ $\times$ $\mathbb{R}$ | p.1 + p.2 * I $\in$ G} $\wedge$
  ComplexAnalysis.R2.TwiceContinuouslyDifferentiable_R2_eps (fun p => (f (p.1 + p.2 * I)).im) {p : $\mathbb{R}$ $\times$ $\mathbb{R}$ | p.1 + p.2 * I $\in$ G}
\end{lstlisting}
    \[ \text{Re}(f), \text{Im}(f) \in C^2(G, \mathbb{R}) \]

    \item \textbf{\texttt{Harmonic\_C\_eps}} ($f : \mathbb{C} \to \mathbb{C}$, $G \subset \mathbb{C}$): $f$ is harmonic on the open set $G$ if its real and imaginary parts satisfy Laplace's equation $\Delta = 0$ on $G$.
\begin{lstlisting}[language=Caml, mathescape=true]
def Harmonic_C_eps (f : $\mathbb{C}$ $\to$ $\mathbb{C}$) (G : Set $\mathbb{C}$) : Prop :=
  ComplexAnalysis.R2.Harmonic_R2_eps (fun p => (f (p.1 + p.2 * I)).re) {p : $\mathbb{R}$ $\times$ $\mathbb{R}$ | p.1 + p.2 * I $\in$ G} $\wedge$
  ComplexAnalysis.R2.Harmonic_R2_eps (fun p => (f (p.1 + p.2 * I)).im) {p : $\mathbb{R}$ $\times$ $\mathbb{R}$ | p.1 + p.2 * I $\in$ G}
\end{lstlisting}
    \[ \Delta \text{Re}(f) = 0, \quad \Delta \text{Im}(f) = 0 \text{ on } G \]

    \item \textbf{\texttt{HarmonicConjugate\_C\_eps}} ($u, v : \mathbb{C} \to \mathbb{R}$, $G \subset \mathbb{C}$): A real-valued function $v$ is a harmonic conjugate of a real-valued function $u$ on an open set $G$ if both are harmonic on $G$ and they satisfy the Cauchy-Riemann equations on $G$.
\begin{lstlisting}[language=Caml, mathescape=true]
def HarmonicConjugate_C_eps (u v : $\mathbb{C}$ $\to$ $\mathbb{R}$) (G : Set $\mathbb{C}$) : Prop :=
  ComplexAnalysis.R2.Harmonic_R2_eps (fun p => u (p.1 + p.2 * I)) {p : $\mathbb{R}$ $\times$ $\mathbb{R}$ | p.1 + p.2 * I $\in$ G} $\wedge$
  ComplexAnalysis.R2.Harmonic_R2_eps (fun p => v (p.1 + p.2 * I)) {p : $\mathbb{R}$ $\times$ $\mathbb{R}$ | p.1 + p.2 * I $\in$ G} $\wedge$
  $\exists$ (ux uy vx vy : $\mathbb{C}$ $\to$ $\mathbb{R}$),
    ($\forall$ z $\in$ G, HasPartialDerivX_C_to_R_eps u (ux z) z) $\wedge$
    ($\forall$ z $\in$ G, HasPartialDerivY_C_to_R_eps u (uy z) z) $\wedge$
    ($\forall$ z $\in$ G, HasPartialDerivX_C_to_R_eps v (vx z) z) $\wedge$
    ($\forall$ z $\in$ G, HasPartialDerivY_C_to_R_eps v (vy z) z) $\wedge$
    ($\forall$ z $\in$ G, ux z = vy z $\wedge$ uy z = -vx z)
\end{lstlisting}
    \[ \Delta u = 0, \quad \Delta v = 0, \quad u_x = v_y, \quad u_y = -v_x \text{ on } G \]
\end{itemize}

\section{\texttt{Sarason/Chapter1.lean}}

\begin{itemize}
    \item \textbf{\texttt{chordalMetric}} ($z, w \in \mathbb{C} \to \mathbb{R}$): The chordal distance between two points in the complex plane mapped to the Riemann sphere.
\begin{lstlisting}[language=Caml, mathescape=true]
noncomputable def chordalMetric (z w : $\mathbb{C}$) : $\mathbb{R}$ := 
  (2 * Complex.abs (z - w)) / (Real.sqrt (1 + Complex.normSq z) * Real.sqrt (1 + Complex.normSq w))
\end{lstlisting}
    \[ \chi(z, w) = \frac{2|z - w|}{\sqrt{1 + |z|^2} \sqrt{1 + |w|^2}} \]

    \item \textbf{\texttt{chordalMetricToInf}} ($z \in \mathbb{C} \to \mathbb{R}$): The chordal distance from a point $z$ to the point at infinity $\infty$.
\begin{lstlisting}[language=Caml, mathescape=true]
noncomputable def chordalMetricToInf (z : $\mathbb{C}$) : $\mathbb{R}$ := 
  2 / Real.sqrt (1 + Complex.normSq z)
\end{lstlisting}
    \[ \chi(z, \infty) = \frac{2}{\sqrt{1 + |z|^2}} \]

    \item \textbf{\texttt{$\Phi$ (Stereographic Projection)}}
\begin{lstlisting}[language=Caml, mathescape=true]
noncomputable def $\Phi$ (z : $\mathbb{C}$) : EuclideanSpace $\mathbb{R}$ (Fin 3) := 
  ![ (2 * z.re) / (1 + Complex.normSq z), 
     (2 * z.im) / (1 + Complex.normSq z), 
     (Complex.normSq z - 1) / (1 + Complex.normSq z) ]
\end{lstlisting}
    \[ \Phi(z) = \left( \frac{2 \operatorname{Re}(z)}{1 + |z|^2}, \frac{2 \operatorname{Im}(z)}{1 + |z|^2}, \frac{|z|^2 - 1}{1 + |z|^2} \right) \in S^2 \]

    \item \textbf{\texttt{ChordalC}}: Type alias for $\mathbb{C}$ equipped with the chordal metric.
\begin{lstlisting}[language=Caml, mathescape=true]
def Chordal$\mathbb{C}$ : Type := $\mathbb{C}$
\end{lstlisting}
    \[ \text{ChordalC} = \mathbb{C} \text{ (equipped with chordal metric)} \]
\end{itemize}

\section{\texttt{Sarason/Chapter2.lean}}

\begin{itemize}
    \item \textbf{\texttt{del} ($\partial$)}
\begin{lstlisting}[language=Caml, mathescape=true]
def del (f : $\mathbb{C}$ $\to$ $\mathbb{C}$) (z : $\mathbb{C}$) : $\mathbb{C}$ :=
  let L := fderiv $\mathbb{R}$ f z
  let df_dx := L 1
  let df_dy := L I
  (1 / 2 : $\mathbb{C}$) * (df_dx - I * df_dy)

/-- The $\partial$$\bar{}$ operator (del-bar) for a function f : $\mathbb{C}$ $\to$ $\mathbb{C}$ at point z. -/
\end{lstlisting}
    \[ \partial f = \frac{1}{2} \left( \frac{\partial f}{\partial x} - i \frac{\partial f}{\partial y} \right) \]

    \item \textbf{\texttt{delBar} ($\bar{\partial}$)}
\begin{lstlisting}[language=Caml, mathescape=true]
def delBar (f : $\mathbb{C}$ $\to$ $\mathbb{C}$) (z : $\mathbb{C}$) : $\mathbb{C}$ :=
  let L := fderiv $\mathbb{R}$ f z
  let df_dx := L 1
  let df_dy := L I
  (1 / 2 : $\mathbb{C}$) * (df_dx + I * df_dy)

-- Local notation for Sarason's flavor
local notation f "$z$" => del f
local notation f "$z$bar" => delBar f

/--
  Relating Sarason's Weierstrass derivative to Mathlib's Filter-based HasDerivAt.
-/
\end{lstlisting}
    \[ \bar{\partial} f = \frac{1}{2} \left( \frac{\partial f}{\partial x} + i \frac{\partial f}{\partial y} \right) \]

    \item \textbf{\texttt{ConformalLinearMap\_eps}} ($ux, uy, vx, vy \in \mathbb{R}$, $v \in \mathbb{C}$): Real-linear transformation on $\mathbb{C} \cong \mathbb{R}^2$ specified by the $2\times 2$ Jacobian entries $(u_x, u_y, v_x, v_y)$.
\begin{lstlisting}[language=Caml, mathescape=true]
def ConformalLinearMap_eps (ux uy vx vy : $\mathbb{R}$) (v : $\mathbb{C}$) : $\mathbb{C}$ :=
  (ux * v.re + uy * v.im) + I * (vx * v.re + vy * v.im)
\end{lstlisting}
    \[ L(v) = (u_x \operatorname{Re}(v) + u_y \operatorname{Im}(v)) + i (v_x \operatorname{Re}(v) + v_y \operatorname{Im}(v)) = (\partial f)\, v + (\bar{\partial} f)\, \bar{v} \]

    \item \textbf{\texttt{del\_eps}} / \textbf{\texttt{delBar\_eps}} ($u_{x0}, u_{y0}, v_{x0}, v_{y0} \in \mathbb{R}$): Wirtinger derivatives expressed purely in terms of the real partial derivatives of $f = u + iv$.
\begin{lstlisting}[language=Caml, mathescape=true]
noncomputable def del_eps (ux uy vx vy : $\mathbb{R}$) : $\mathbb{C}$ :=
  (1 / 2 : $\mathbb{C}$) * ( (ux + vy) + I * (vx - uy) )
\end{lstlisting}
    \[ \partial f = \frac{1}{2} (u_{x0} + v_{y0}) + \frac{i}{2} (v_{x0} - u_{y0}) \]

    \item \textbf{\texttt{delBar\_eps}}
\begin{lstlisting}[language=Caml, mathescape=true]
noncomputable def delBar_eps (ux uy vx vy : $\mathbb{R}$) : $\mathbb{C}$ :=
  (1 / 2 : $\mathbb{C}$) * ( (ux - vy) + I * (vx + uy) )
\end{lstlisting}
    \[ \bar{\partial} f = \frac{1}{2} (u_{x0} - v_{y0}) + \frac{i}{2} (v_{x0} + u_{y0}) \]

    \item \textbf{\texttt{path\_in\_C}} ($x, y \in \mathbb{C}$): A continuous path $\gamma : [0, 1] \to \mathbb{C}$ starting at $x$ and ending at $y$.
\begin{lstlisting}[language=Caml, mathescape=true]
def path_in_C (x y : $\mathbb{C}$) : Type := Path x y

/-- Velocity / derivative of a path $\gamma$ : path_in_C x y at parameter time t $\in$ $\mathbb{R}$. -/
\end{lstlisting}
    \[ \gamma : [0, 1] \to \mathbb{C} \text{ where } \gamma(0)=x, \gamma(1)=y \]

    \item \textbf{\texttt{translated\_path}} ($c, x, y \in \mathbb{C}$, $\gamma : \text{path\_in\_C}\ x\ y$): A path shifted by the constant $c$, creating a path from $c+x$ to $c+y$.
\begin{lstlisting}[language=Caml, mathescape=true]
noncomputable def translated_path {x y : $\mathbb{C}$} (c : $\mathbb{C}$) ($\gamma$ : path_in_C x y) : path_in_C (c + x) (c + y) where
  toFun := fun t => c + $\gamma$.toFun t
  continuous_toFun := Continuous.add continuous_const $\gamma$.continuous_toFun
  source' := by { change c + $\gamma$.toFun 0 = c + x; have h : $\gamma$.toFun 0 = x := ($\gamma$ : Path x y).source'; rw [h] }
  target' := by { change c + $\gamma$.toFun 1 = c + y; have h : $\gamma$.toFun 1 = y := ($\gamma$ : Path x y).target'; rw [h] }
\end{lstlisting}
    \[ \gamma_{c}(t) = c + \gamma(t) \]

    \item \textbf{\texttt{pathDeriv}} ($\gamma : \text{path\_in\_C}\ x\ y$, $t \in \mathbb{R}$): The tangent vector $\gamma'(t) \in \mathbb{C}$, extracted via choice from \texttt{HasDerivAt\_R\_to\_C\_eps} extended to $\mathbb{R}$.
\begin{lstlisting}[language=Caml, mathescape=true]
noncomputable def pathDeriv {x y : $\mathbb{C}$} ($\gamma$ : path_in_C x y) (t : $\mathbb{R}$) : $\mathbb{C}$ :=
  deriv_R_to_C_eps $\gamma$.extend t

/-- The direction of a path $\gamma$ at time t is the argument of its derivative. -/
\end{lstlisting}
    \[ \gamma'(t) = \frac{d\gamma}{dt} \]

    \item \textbf{\texttt{pathDirection}} ($\gamma : \text{path\_in\_C}\ x\ y$, $t \in \mathbb{R}$): The argument (angle) of the tangent vector, $\arg(\gamma'(t))$.
\begin{lstlisting}[language=Caml, mathescape=true]
noncomputable def pathDirection {x y : $\mathbb{C}$} ($\gamma$ : path_in_C x y) (t : $\mathbb{R}$) : $\mathbb{R}$ :=
  arg (pathDeriv $\gamma$ t)

/-- The angle between two paths $\gamma$$_1$ and $\gamma$$_2$ at times t$_1$ and t$_2$ is the difference of their directions. -/
\end{lstlisting}
    \[ \text{Direction}(\gamma, t) = \arg(\gamma'(t)) \]

    \item \textbf{\texttt{pathAngle}} ($\gamma_1, t_1, \gamma_2, t_2$): The angle between two intersecting paths.
\begin{lstlisting}[language=Caml, mathescape=true]
noncomputable def pathAngle {x$_1$ y$_1$ x$_2$ y$_2$ : $\mathbb{C}$} ($\gamma$$_1$ : path_in_C x$_1$ y$_1$) (t$_1$ : $\mathbb{R}$) ($\gamma$$_2$ : path_in_C x$_2$ y$_2$) (t$_2$ : $\mathbb{R}$) : $\mathbb{R}$ :=
  pathDirection $\gamma$$_2$ t$_2$ - pathDirection $\gamma$$_1$ t$_1$

/--
  The angle between two paths is equal to the argument of the product of the second path's derivative
  and the conjugate of the first path's derivative, modulo 2$\pi$.
  We prove this equality here assuming the angle difference stays within the principal branch bounds.
-/
\end{lstlisting}
    \[ \text{Angle}(\gamma_1, t_1, \gamma_2, t_2) = \arg(\gamma_2'(t_2)) - \arg(\gamma_1'(t_1)) \]

    \item \textbf{\texttt{conformal}} ($f : \mathbb{C} \to \mathbb{C}$, $z_0 \in \mathbb{C}$): $f$ preserves angles between any two paths intersecting at $z_0$ with non-zero derivatives.
\begin{lstlisting}[language=Caml, mathescape=true]
def conformal (f : $\mathbb{C}$ $\to$ $\mathbb{C}$) (z$_0$ : $\mathbb{C}$) : Prop :=
  $\forall$ (x$_1$ y$_1$ x$_2$ y$_2$ : $\mathbb{C}$) ($\gamma$$_1$ : path_in_C x$_1$ y$_1$) (t$_1$ : $\mathbb{R}$) ($\gamma$$_2$ : path_in_C x$_2$ y$_2$) (t$_2$ : $\mathbb{R}$) (d$_1$ d$_2$ : $\mathbb{C}$),
    $\gamma$$_1$.extend t$_1$ = z$_0$ $\to$
    $\gamma$$_2$.extend t$_2$ = z$_0$ $\to$
    HasDerivAt_R_to_C_eps $\gamma$$_1$.extend (pathDeriv $\gamma$$_1$ t$_1$) t$_1$ $\to$
    HasDerivAt_R_to_C_eps $\gamma$$_2$.extend (pathDeriv $\gamma$$_2$ t$_2$) t$_2$ $\to$
    pathDeriv $\gamma$$_1$ t$_1$ $\neq$ 0 $\to$
    pathDeriv $\gamma$$_2$ t$_2$ $\neq$ 0 $\to$
    -- HasDerivAt_R_to_C_eps f f' t says f' is the derivative of f at t.
    HasDerivAt_R_to_C_eps (f $\circ$ $\gamma$$_1$.extend) d$_1$ t$_1$ $\to$
    HasDerivAt_R_to_C_eps (f $\circ$ $\gamma$$_2$.extend) d$_2$ t$_2$ $\to$
    d$_1$ $\neq$ 0 $\wedge$ d$_2$ $\neq$ 0 $\wedge$ arg d$_2$ - arg d$_1$ = pathAngle $\gamma$$_1$ t$_1$ $\gamma$$_2$ t$_2$

/--
  \SII.11 Angle preservation at a point.  CONFORMALITY.
  If f is holomorphic at z$_0$ with non-zero derivative, then it preserves angles between regular paths intersecting at z$_0$.
-/
\end{lstlisting}
    \[ \text{Conformal } f \text{ at } z_0 \iff f \text{ preserves angles between all regular intersecting curves at } z_0 \]

    \item \textbf{\texttt{line\_path\_zero}} ($v \in \mathbb{C}$): A straight-line path from $-v$ to $v$ passing through $0$ at $t = 1/2$.
\begin{lstlisting}[language=Caml, mathescape=true]
noncomputable def line_path_zero (v : $\mathbb{C}$) : path_in_C (-v) v where
  toFun := fun t => -v + (2 * ((t : $\mathbb{R}$) : $\mathbb{C}$)) * v
\end{lstlisting}
    \[ \gamma(t) = -v + 2tv = (2t - 1)v \]

    \item \textbf{\texttt{centered\_line\_path}} ($z, v \in \mathbb{C}$): A straight-line path from $z - v$ to $z + v$ evaluating to $z$ at $t = 1/2$.
\begin{lstlisting}[language=Caml, mathescape=true]
noncomputable def centered_line_path (z v : $\mathbb{C}$) : path_in_C (z - v) (z + v) where
  toFun := fun t => z + (2 * (t : $\mathbb{R}$) - 1) * v
\end{lstlisting}
    \[ \gamma(t) = z + (2t - 1)v \]
\end{itemize}

\section{\texttt{Sarason/Chapter3.lean}}

\begin{itemize}
    \item \textbf{\texttt{linFracTrans}} ($a, b, c, d, z \in \mathbb{C}$): A linear-fractional transformation (M\"obius transformation) on $\mathbb{C}$.
\begin{lstlisting}[language=Caml, mathescape=true]
noncomputable def linFracTrans (a b c d z : $\mathbb{C}$) : $\mathbb{C}$ :=
  (a * z + b) / (c * z + d)
\end{lstlisting}
    \[ \varphi(z) = \frac{az + b}{cz + d} \]

    \item \textbf{\texttt{$\varphi$}} ($a, b, c, d, z \in \mathbb{C}$): Sarason's shorthand notation for \texttt{linFracTrans}.
\begin{lstlisting}[language=Caml, mathescape=true]
noncomputable abbrev $\varphi$ (a b c d z : $\mathbb{C}$) : $\mathbb{C}$ :=
  linFracTrans a b c d z
\end{lstlisting}
    \[ \varphi(a, b, c, d, z) = \frac{az + b}{cz + d} \]

    \item \textbf{\texttt{linFracTransExt}} ($a, b, c, d \in \mathbb{C}$, $w \in \mathbb{C} \cup \{\infty\}$): The extended linear-fractional transformation on the Riemann sphere $\mathbb{C} \cup \{\infty\}$ (\texttt{OnePoint} $\mathbb{C}$).
\begin{lstlisting}[language=Caml, mathescape=true]
noncomputable def linFracTransExt (a b c d : $\mathbb{C}$) : OnePoint $\mathbb{C}$ $\to$ OnePoint $\mathbb{C}$
  | (z : $\mathbb{C}$) => if c * z + d = 0 then $\infty$ else (linFracTrans a b c d z : $\mathbb{C}$)
  | $\infty$ => if c = 0 then $\infty$ else (a / c : $\mathbb{C}$)
\end{lstlisting}
    \[ \varphi(w) = \begin{cases} \frac{az+b}{cz+d} & \text{if } w = z \in \mathbb{C} \text{ and } cz+d \neq 0, \\ \infty & \text{if } w = z \in \mathbb{C} \text{ and } cz+d = 0, \\ \infty & \text{if } w = \infty \text{ and } c = 0, \\ \frac{a}{c} & \text{if } w = \infty \text{ and } c \neq 0. \end{cases} \]

    \item \textbf{\texttt{IsLinFracTrans}} ($\varphi : \mathbb{C} \cup \{\infty\} \to \mathbb{C} \cup \{\infty\}$): Predicate stating that $\varphi$ is a non-degenerate linear-fractional transformation on $\mathbb{C} \cup \{\infty\}$.
\begin{lstlisting}[language=Caml, mathescape=true]
def IsLinFracTrans ($\varphi$ : OnePoint $\mathbb{C}$ $\to$ OnePoint $\mathbb{C}$) : Prop :=
  $\exists$ a b c d : $\mathbb{C}$, a * d - b * c $\neq$ 0 $\wedge$ $\varphi$ = linFracTransExt a b c d
\end{lstlisting}
    \[ \operatorname{IsLinFracTrans}(\varphi) \iff \exists a, b, c, d \in \mathbb{C},\; ad - bc \neq 0 \;\wedge\; \varphi(w) = \frac{aw + b}{cw + d} \]

    \item \textbf{\texttt{gl2OfCoeffs}} ($a, b, c, d \in \mathbb{C}$, $ad - bc \neq 0$): The invertible $2 \times 2$ matrix in $\operatorname{GL}(2, \mathbb{C})$ associated with coefficients $(a, b, c, d)$.
\begin{lstlisting}[language=Caml, mathescape=true]
noncomputable def gl2OfCoeffs (a b c d : $\mathbb{C}$) (h_det : a * d - b * c $\neq$ 0) : GL (Fin 2) $\mathbb{C}$ :=
  Matrix.GeneralLinearGroup.mkOfDetNeZero !![a, b; c, d] ...
\end{lstlisting}
    \[ \begin{pmatrix} a & b \\ c & d \end{pmatrix} \in \operatorname{GL}(2, \mathbb{C}), \quad ad - bc \neq 0 \]

    \item \textbf{\texttt{translationMatrix}} ($b \in \mathbb{C}$): 
\begin{lstlisting}[language=Caml, mathescape=true]
noncomputable def translationMatrix (b : $\mathbb{C}$) : GL (Fin 2) $\mathbb{C}$ := 
  Matrix.GeneralLinearGroup.mkOfDetNeZero !![1, b; 0, 1] ...
\end{lstlisting}
    \[ \begin{pmatrix} 1 & b \\ 0 & 1 \end{pmatrix} \]

    \item \textbf{\texttt{scalingMatrix}} ($a \in \mathbb{C}, a \neq 0$): 
\begin{lstlisting}[language=Caml, mathescape=true]
noncomputable def scalingMatrix (a : $\mathbb{C}$) (ha : a $\neq$ 0) : GL (Fin 2) $\mathbb{C}$ := 
  Matrix.GeneralLinearGroup.mkOfDetNeZero !![a, 0; 0, 1] ...
\end{lstlisting}
    \[ \begin{pmatrix} a & 0 \\ 0 & 1 \end{pmatrix} \]

    \item \textbf{\texttt{inversionMatrix}}: 
\begin{lstlisting}[language=Caml, mathescape=true]
noncomputable def inversionMatrix : GL (Fin 2) $\mathbb{C}$ := 
  Matrix.GeneralLinearGroup.mkOfDetNeZero !![0, 1; 1, 0] ...
\end{lstlisting}
    \[ \begin{pmatrix} 0 & 1 \\ 1 & 0 \end{pmatrix} \]
\end{itemize}

\end{document}